% GENERATED FILE. Edit canonical lessons, then run npm run generate:books.
% canonical-source-hash: fnv1a64:7a6985f138da4ae5
% canonical-lessons: IT-C301-soccorso, IT-C301-soccorrere, IT-C301-soccorritore, IT-C301-barella, IT-C301-sirena

\chapter{Arrivano i soccorsi}
\label{ch:it-arrivano-i-soccorsi}
\hlchaptermodality{301}

\begin{chapteropening}
\textbf{By the end of this chapter:} \emph{I can describe the rescue: help arriving, the rescuers, the stretcher and the siren.}

Complete the last lesson of chapter 301: I can describe the rescue: help arriving, the rescuers, the stretcher and the siren.
\end{chapteropening}

\section[il soccorso]{il soccorso — emergency help, rescue}

\label{lesson:IT-C301-soccorso}



\begin{quote}
Before the new one: say the Italian for unhurt, then the Italian for to come through all right.
\end{quote}

\subsection*{You'll want to know: il soccorso}
\textbf{il soccorso} — \textquotedblleft{}emergency help, rescue\textquotedblright{}.

\begin{quote}
Il soccorso è arrivato dopo dieci minuti.
\end{quote}

\begin{quote}
\emph{Help arrived after ten minutes.}
\end{quote}

\emph{Chiamare i soccorsi} — to call for help.

\begin{cousinweb}
\emph{soccorso} is inherited from Vulgar Latin \emph{succursum} (EtymDB), built on Latin \emph{succurrere} \textquotedblleft{}to run to help\textquotedblright{}.
\end{cousinweb}

\begin{grammarlens}[title={What you've built}]
One more word for telling, explaining and describing at B1.
\end{grammarlens}

\subsection*{Your turn}
\begin{itemize}
\raggedright
  \item \emph{Say it:} \emph{il soccorso}
  \item \emph{Say it:} \emph{il soccorso}, once more
  \item \emph{Say it:} a sentence of your own with \emph{il soccorso}
  \item \emph{Recall:} say the Italian for unhurt, then the Italian for to come through all right, then say \emph{il soccorso} again
\end{itemize}

\begin{tcolorbox}[breakable,colback=teal!4,colframe=teal!35!black,title={Before you move on}]
What does il soccorso mean? (\textquotedblleft{}Emergency help, rescue\textquotedblright{}.) Say it once more.
\end{tcolorbox}
\section[soccorrere]{soccorrere — to come to the aid of}

\label{lesson:IT-C301-soccorrere}



\begin{quote}
Before the new one: say the Italian for to come through all right, then the Italian for emergency help.
\end{quote}

\subsection*{You'll want to know: soccorrere}
\textbf{soccorrere} — \textquotedblleft{}to come to the aid of\textquotedblright{}.

\begin{quote}
I passanti hanno soccorso il ciclista.
\end{quote}

\begin{quote}
\emph{Passers-by came to the cyclist's aid.}
\end{quote}

Conjugated like \emph{correre}: \emph{ha soccorso}.

\begin{cousinweb}
\emph{soccorrere} is inherited from Latin \emph{succurrere} \textquotedblleft{}to run up to help\textquotedblright{} (EtymDB), from \emph{sub} and \emph{currere} \textquotedblleft{}to run\textquotedblright{}.
\end{cousinweb}

\begin{grammarlens}[title={What you've built}]
One more word for telling, explaining and describing at B1.
\end{grammarlens}

\subsection*{Your turn}
\begin{itemize}
\raggedright
  \item \emph{Say it:} \emph{soccorrere}
  \item \emph{Say it:} \emph{soccorrere}, once more
  \item \emph{Say it:} a sentence of your own with \emph{soccorrere}
  \item \emph{Recall:} say the Italian for to come through all right, then the Italian for emergency help, then say \emph{soccorrere} again
\end{itemize}

\begin{tcolorbox}[breakable,colback=teal!4,colframe=teal!35!black,title={Before you move on}]
What does soccorrere mean? (\textquotedblleft{}To come to the aid of\textquotedblright{}.) Say it once more.
\end{tcolorbox}
\section[il soccorritore]{il soccorritore — the rescuer, the paramedic}

\label{lesson:IT-C301-soccorritore}



\begin{quote}
Before the new one: say the Italian for emergency help, then the Italian for to come to the aid of.
\end{quote}

\subsection*{You'll want to know: il soccorritore}
\textbf{il soccorritore} — \textquotedblleft{}the rescuer, the paramedic\textquotedblright{}.

\begin{quote}
I soccorritori hanno lavorato tutta la notte.
\end{quote}

\begin{quote}
\emph{The rescuers worked all night.}
\end{quote}

A woman is \emph{la soccorritrice}.

\begin{cousinweb}
\emph{soccorritore} is \emph{soccorrere} with the agent ending \emph{-tore}; the root is Latin \emph{currere} \textquotedblleft{}to run\textquotedblright{} (EtymDB).
\end{cousinweb}

\begin{grammarlens}[title={What you've built}]
One more word for telling, explaining and describing at B1.
\end{grammarlens}

\subsection*{Your turn}
\begin{itemize}
\raggedright
  \item \emph{Say it:} \emph{il soccorritore}
  \item \emph{Say it:} \emph{il soccorritore}, once more
  \item \emph{Say it:} a sentence of your own with \emph{il soccorritore}
  \item \emph{Recall:} say the Italian for emergency help, then the Italian for to come to the aid of, then say \emph{il soccorritore} again
\end{itemize}

\begin{tcolorbox}[breakable,colback=teal!4,colframe=teal!35!black,title={Before you move on}]
What does il soccorritore mean? (\textquotedblleft{}The rescuer, the paramedic\textquotedblright{}.) Say it once more.
\end{tcolorbox}
\section[la barella]{la barella — the stretcher}

\label{lesson:IT-C301-barella}



\begin{quote}
Before the new one: say the Italian for to come to the aid of, then the Italian for the rescuer.
\end{quote}

\subsection*{You'll want to know: la barella}
\textbf{la barella} — \textquotedblleft{}the stretcher\textquotedblright{}.

\begin{quote}
Lo hanno messo sulla barella.
\end{quote}

\begin{quote}
\emph{They put him on the stretcher.}
\end{quote}

\emph{In barella} — on a stretcher.

\begin{cousinweb}
The origin of \emph{barella} is not settled; it is often linked with \emph{bara} \textquotedblleft{}a bier\textquotedblright{}, which EtymDB traces to a Longobard word.
\end{cousinweb}

\begin{grammarlens}[title={What you've built}]
One more word for telling, explaining and describing at B1.
\end{grammarlens}

\subsection*{Your turn}
\begin{itemize}
\raggedright
  \item \emph{Say it:} \emph{la barella}
  \item \emph{Say it:} \emph{la barella}, once more
  \item \emph{Say it:} a sentence of your own with \emph{la barella}
  \item \emph{Recall:} say the Italian for to come to the aid of, then the Italian for the rescuer, then say \emph{la barella} again
\end{itemize}

\begin{tcolorbox}[breakable,colback=teal!4,colframe=teal!35!black,title={Before you move on}]
What does la barella mean? (\textquotedblleft{}The stretcher\textquotedblright{}.) Say it once more.
\end{tcolorbox}
\section[la sirena]{la sirena — the siren}

\label{lesson:IT-C301-sirena}



\begin{quote}
Before the new one: say the Italian for the rescuer, then the Italian for the stretcher.
\end{quote}

\subsection*{You'll want to know: la sirena}
\textbf{la sirena} — \textquotedblleft{}the siren\textquotedblright{}.

\begin{quote}
Si sente la sirena dell'ambulanza.
\end{quote}

\begin{quote}
\emph{You can hear the ambulance siren.}
\end{quote}

In stories, \emph{la sirena} is also a mermaid.

\begin{cousinweb}
\emph{sirena} comes through Late Latin \emph{sīrēna} from Greek \emph{Seir\'{\={e}}n} (EtymDB), the singing sea creatures of the old myths.
\end{cousinweb}

\begin{grammarlens}[title={What you've built}]
One more word for telling, explaining and describing at B1.
\end{grammarlens}

\subsection*{Your turn}
\begin{itemize}
\raggedright
  \item \emph{Say it:} \emph{la sirena}
  \item \emph{Say it:} \emph{la sirena}, once more
  \item \emph{Say it:} a sentence of your own with \emph{la sirena}
  \item \emph{Recall:} say the Italian for the rescuer, then the Italian for the stretcher, then say \emph{la sirena} again
\end{itemize}

\begin{tcolorbox}[breakable,colback=teal!4,colframe=teal!35!black,title={Before you move on}]
What does la sirena mean? (\textquotedblleft{}The siren\textquotedblright{}.) Say it once more.
\end{tcolorbox}
